\chapter{Modelo MVC y API REST}
\label{cap:mvc-api-rest}

\section{Flujo MVC de una petici\'on autenticada}
\label{sec:flujo-mvc}

BIOPET implementa el patr\'on Modelo-Vista-Controlador de forma est\'andar
para una aplicaci\'on Spring Boot: el cliente Angular consume la API REST
del backend, que separa controladores (\texttt{controller}), l\'ogica de
negocio (\texttt{service}) y acceso a datos (\texttt{repository}/
\texttt{entity}). El README ra\'iz del repositorio documenta, paso a paso
y verificado contra el c\'odigo, el ciclo de vida completo de una
petici\'on autenticada (\texttt{GET /api/mascotas}):

\begin{enumerate}[nosep]
  \item \textbf{Cliente Angular:} \texttt{MascotaApiService.listar()} env\'ia
    la petici\'on incluyendo la cookie de autenticaci\'on.
  \item \textbf{JwtAuthenticationFilter:} intercepta la petici\'on y
    resuelve el JWT desde la cookie.
  \item \textbf{SecurityContextHolder:} Spring registra la autenticaci\'on
    del usuario ya validado.
  \item \textbf{MascotaController:} recibe la petici\'on autenticada y los
    par\'ametros de paginaci\'on.
  \item \textbf{MascotaService:} ejecuta la l\'ogica de negocio (incluye la
    regla de acceso por rol y por propiedad) dentro de una transacci\'on de
    solo lectura.
  \item \textbf{UsuarioRepository / MascotaRepository:} obtienen al
    usuario autenticado y, seg\'un su rol, ejecutan la consulta
    correspondiente (global o filtrada por due\~no).
  \item \textbf{PostgreSQL:} Hibernate genera y ejecuta la consulta SQL.
  \item \textbf{Jackson:} serializa la respuesta a JSON.
  \item \textbf{Cliente Angular:} recibe la respuesta \texttt{200 OK} con
    el contenido paginado.
\end{enumerate}

\evidencia{../../diagramas/flujo-mvc-springboot.png}{Diagrama de secuencia del flujo MVC de BIOPET (cliente Angular a PostgreSQL).}{fig:flujo-mvc}

Este mismo patr\'on de capas (\texttt{Controller} $\to$ \texttt{Service}
$\to$ \texttt{Repository}) se repite consistentemente en los dem\'as
m\'odulos descritos en la secci\'on~\ref{sec:proposito-alcance} (Citas,
Consultas, Vacunas, Usuarios), verificado directamente en el c\'odigo
fuente del backend.

\section{Endpoints de la API REST}
\label{sec:endpoints-rest}

BIOPET expone una API REST organizada por recurso. La siguiente tabla
resume los grupos de endpoints reales, verificados contra los controllers
actuales del backend (\texttt{Backend/src/main/java/com/biopet/controller/}):

\begin{table}[H]
  \centering
  \small
  \caption{Recursos REST de BIOPET, verificados contra el c\'odigo actual.}
  \label{tab:recursos-rest}
  \begin{tabular}{@{}lll@{}}
    \toprule
    Recurso & Ruta base & Operaciones \\
    \midrule
    Autenticaci\'on & \texttt{/api/auth} & registro, login, refresh, logout \\
    Usuarios & \texttt{/api/usuarios} & \texttt{/me} + CRUD administrativo \\
    Mascotas & \texttt{/api/mascotas} & CRUD + resumen por especie \\
    Citas & \texttt{/api/citas} & CRUD \\
    Consultas & \texttt{/api/consultas} & CRUD \\
    Vacunas & \texttt{/api/vacunas} & CRUD + listado por mascota \\
    API externa & \texttt{/api/externa/especies} & consulta (con cach\'e Redis) \\
    \bottomrule
  \end{tabular}
\end{table}

Todas las respuestas de error de la API usan el formato uniforme
\texttt{ProblemDetail} (RFC 7807)~\parencite{rfc7807}, y todos los recursos
protegidos requieren la cookie de sesi\'on emitida por
\texttt{/api/auth/login}, con autorizaci\'on adicional por rol
(\texttt{@PreAuthorize}) y, en varios recursos, por propiedad del registro.

\section{Documentaci\'on OpenAPI / Swagger}
\label{sec:openapi-swagger}

BIOPET documenta su API mediante \texttt{springdoc-openapi}, siguiendo la
especificaci\'on OpenAPI~\parencite{openapi-spec}. La configuraci\'on real
(\url{Backend/src/main/resources/application.yml}) expone:

\begin{itemize}[nosep]
  \item \textbf{Swagger UI:} \texttt{/api/docs}.
  \item \textbf{Documento OpenAPI (JSON):} \texttt{/api/openapi}.
\end{itemize}

Ambas rutas son p\'ublicas (no requieren autenticaci\'on), verificado en
\texttt{SecurityConfig} y probado autom\'aticamente por
\texttt{SwaggerUiTest}: la primera prueba confirma que \url{/api/docs}
redirige a la interfaz real de Swagger UI (cuyo cuerpo contiene el
marcador \texttt{swagger-ui}); la segunda confirma que \url{/api/openapi}
responde \texttt{200} con un documento OpenAPI v\'alido cuyo t\'itulo es
\texttt{BIOPET API}.

\evidencia{../evidencias/jaime/swagger-ui.png}{Swagger UI de BIOPET, accesible en /api/docs, mostrando los grupos de endpoints reales.}{fig:swagger-ui}

\section{Evidencia Postman / Newman}
\label{sec:evidencia-postman}

El repositorio incluye colecciones Postman reales y reproducibles,
verificadas contra el c\'odigo actual y ejecutadas de punta a punta con
Newman:

\begin{itemize}[nosep]
  \item \texttt{docs/postman/BIOPET.postman\_collection.json} --- colecci\'on
    principal (Auth + Mascotas): 44 requests definidos en 7 carpetas.
    Ejecuci\'on real con Newman: 46 requests, 234 assertions, 0 fallos.
  \item \texttt{docs/postman/BIOPET-Vacunas.postman\_collection.json} ---
    colecci\'on adicional del CRUD de Vacunas: 12 requests.
\end{itemize}

Esto demuestra, muy por encima del m\'inimo de tres endpoints exigido por
la actividad, que la API REST de BIOPET funciona de extremo a extremo
contra el backend real, incluyendo casos de \'exito, de validaci\'on
(\texttt{422}), de autenticaci\'on (\texttt{401}), de autorizaci\'on
(\texttt{403}) y de limitaci\'on de intentos de login (\texttt{429}).
